\chapter{Ruang Vektor Bernorma}\label{ch:b2:nvs}

Ketika \omterm{def:b2:metric:def}{ruang metriknya} berupa ruang vektor dan jaraknya berasal dari
sebuah \omterm{def:b2:nvs:norm}{norma}, topologi dan aljabar linear mulai saling berjalin: pemetaan
linear bersifat \omterm{def:b2:metric:continuity}{kontinu} persis ketika ia terbatas pada bola satuan,
dimensi hingga memaksa semua \omterm{def:b2:nvs:norm}{norma} sepakat, dan kelengkapan mengubah
deret yang konvergen mutlak menjadi deret konvergen. Jurang antara
dimensi hingga dan tak hingga --- yang mengkristal pada teorema Riesz
--- adalah pelajaran terdalam bab ini.

Di sepanjang bab ini, $E, F$ adalah ruang vektor atas $K = \R$ atau $\C$.

\section{Norma}

\begin{definition}\label{def:b2:nvs:norm}
Sebuah \emph{norma}\index{norma} pada $E$ adalah pemetaan
$\norm{\,\cdot\,} \colon E \to \R_+$ yang, untuk setiap $x, y \in E$ dan
$\lambda \in K$, memenuhi
\[
\norm x = 0 \iff x = 0,
\qquad
\norm{\lambda x} = \abs\lambda\,\norm x,
\qquad
\norm{x + y} \leq \norm x + \norm y .
\]
Maka $d(x, y) = \norm{x - y}$ merupakan jarak, dan seluruh
\cref{ch:b2:metric} berlaku. Ketaksamaan segitiga terbalik
$\bigl|\norm x - \norm y\bigr| \leq \norm{x - y}$ membuat normanya
sendiri \omterm{def:b2:metric:continuity}{Lipschitz} berkonstanta $1$; sedangkan penjumlahan dan perkalian
skalar bersifat \omterm{def:b2:metric:continuity}{kontinu} (lewat taksiran $\norm{(x + y) - (x' + y')} \leq
\norm{x - x'} + \norm{y - y'}$, dan seterusnya).
\end{definition}

\begin{example}\label{ex:b2:nvs:examples}
Pada $K^n$:
\[
\norm{x}_1 = \sum_i \abs{x_i},
\qquad
\norm{x}_2 = \Bigl(\sum_i \abs{x_i}^2\Bigr)^{1/2},
\qquad
\norm{x}_\infty = \max_i \abs{x_i}
\]
($\norm\cdot_2$ merupakan \omterm{def:b2:nvs:norm}{norma} berkat Cauchy--Schwarz, jilid Tahun
ke-1). Pada $C(\intcc{a}{b})$:
\[
\norm f_\infty = \sup \abs f,
\qquad
\norm f_1 = \int_a^b \abs f,
\qquad
\norm f_2 = \Bigl(\int_a^b \abs f^2\Bigr)^{1/2},
\]
dua yang terakhir menjadi \omterm{def:b2:nvs:norm}{norma} berkat kepositifan tegas
integralnya dan Cauchy--Schwarz untuk integral (jilid Tahun ke-1). Pada
matriks: \omterm{def:b2:nvs:norm}{norma} apa pun pada $\mathcal M_n(K) \simeq K^{n^2}$; adapun
\omterm{thm:b2:nvs:continuouslinear}{norma operator} di bawah ini yang penting secara struktural.
\end{example}

\begin{definition}[Norma yang setara]\label{def:b2:nvs:equivalent}
Dua \omterm{def:b2:nvs:norm}{norma} $N_1, N_2$ pada $E$ disebut \emph{setara}\index{norma
setara} apabila ada konstanta $c, C > 0$ dengan
\[
c\,N_1 \leq N_2 \leq C\, N_1 .
\]
\omterm{def:b2:nvs:norm}{Norma} yang setara punya himpunan \omterm{def:b2:metric:topology}{terbuka} yang sama, barisan konvergen dan
barisan Cauchy yang sama, himpunan bagian \omterm{def:b2:metric:compact}{kompak} dan \omterm{def:b2:metric:complete}{lengkap} yang sama:
jadi analisis yang sama.
\end{definition}

\begin{example}[Ketaksetaraan dalam dimensi tak hingga]\label{ex:b2:nvs:nonequivalent}
Pada $C(\intcc{0}{1})$: selalu berlaku $\norm f_1 \leq \norm f_\infty$,
tetapi tak ada batas sebaliknya: sebab $f_n(x) = x^n$ punya
$\norm{f_n}_\infty = 1$ dan $\norm{f_n}_1 = \frac{1}{n+1} \to 0$. Jadi
$f_n \to 0$ untuk $\norm\cdot_1$ tetapi tidak untuk $\norm\cdot_\infty$:
kedua \omterm{def:b2:nvs:norm}{norma} itu berselisih justru tentang kekonvergenannya sendiri.
\end{example}

\begin{example}[Konstanta yang gamblang dalam dimensi $n$]\label{ex:b2:nvs:constants}
Pada $K^n$ ketiga \omterm{def:b2:nvs:norm}{norma} klasiknya \omterm{def:b2:nvs:equivalent}{setara} dengan konstanta
yang tajam:
\[
\norm x_\infty \leq \norm x_2 \leq \norm x_1
\leq \sqrt n\,\norm x_2 \leq n\,\norm x_\infty ,
\]
dengan batas tengahnya $\norm x_1 \leq \sqrt n\norm x_2$ berasal dari
Cauchy--Schwarz terhadap vektor yang semua unsurnya satu. Vektor
ekstremalnya: $e_1$ membuat dua ketaksamaan pertamanya menjadi kesamaan,
dan $(1, 1, \dots, 1)$ membuat dua yang terakhir demikian. Dimensi $n$
duduk kasatmata di dalam konstantanya --- inilah benih kuantitatif
kegagalannya dalam dimensi tak hingga: ketika $n \to \infty$ tak ada
konstanta seragam yang bertahan, dan persis itulah yang diperlihatkan
\cref{ex:b2:nvs:nonequivalent} pada ruang fungsi.
\end{example}

\begin{omfigure}
\begin{tikzpicture}[scale=1.5]
  \draw[->] (-1.45,0) -- (1.5,0) node[below, font=\small] {$x$};
  \draw[->] (0,-1.45) -- (0,1.5) node[left, font=\small] {$y$};
  \draw[omProp, very thick]
    (1,1) -- (-1,1) -- (-1,-1) -- (1,-1) -- cycle;
  \draw[omThm, very thick] (0,0) circle (1);
  \draw[omDef, very thick]
    (1,0) -- (0,1) -- (-1,0) -- (0,-1) -- cycle;
  \node[omDef, font=\small] at (0.36,0.28)
    {$\norm\cdot_1$};
  \node[omThm, font=\small] at (0.55,0.92)
    {$\norm\cdot_2$};
  \node[omProp, font=\small] at (1.24,1.13)
    {$\norm\cdot_\infty$};
\end{tikzpicture}

{\small Bola satuan ketiga \omterm{def:b2:nvs:norm}{norma} klasik $\R^2$,
yang bersarang persis seperti dituntut ketaksamaan pada
\cref{ex:b2:nvs:constants}: makin kecil bolanya, makin besar \omterm{def:b2:nvs:norm}{normanya}.
Kebulatannya penting: sisi datar pada belah ketupat dan pada perseginya
tak lain kegagalan kecembungan tegas yang dimanfaatkan soal akhir pekan
\cref{ch:b2:realfun} dan soal akhir pekan bab ini (pertanyaan 4).}
\end{omfigure}

\section{Pemetaan linear yang kontinu}

\begin{theorem}[Pencirian]\label{thm:b2:nvs:continuouslinear}
Untuk pemetaan linear $u \colon E \to F$ antara ruang bernorma,
pernyataan berikut \omterm{def:b2:nvs:equivalent}{setara}:
\begin{enumerate}
  \item $u$ \omterm{def:b2:metric:continuity}{kontinu};
  \item $u$ \omterm{def:b2:metric:continuity}{kontinu} di $0$;
  \item $u$ terbatas pada bola satuan tertutup: $\sup_{\norm x \leq
        1} \norm{u(x)} < \infty$;
  \item ada $C \geq 0$ dengan $\norm{u(x)} \leq C \norm x$ untuk
        setiap $x$;
  \item $u$ bersifat \omterm{def:b2:metric:continuity}{Lipschitz}.
\end{enumerate}
Nilai $C$ terkecil yang demikian disebut \emph{norma
operator}\index{norma operator} $\vertiii{u} = \sup_{\norm x \leq
1}\norm{u(x)} = \sup_{x \neq 0} \frac{\norm{u(x)}}{\norm x}$; \omterm{def:b2:nvs:norm}{norma} itu
menjadikan ruang $\mathcal{L}_c(E, F)$ berisi pemetaan linear \omterm{def:b2:metric:continuity}{kontinu}
sebuah ruang bernorma, dengan
\[
\vertiii{v \circ u} \leq \vertiii v\, \vertiii u .
\]
\end{theorem}

\begin{proof}
(1 $\Rightarrow$ 2) sepele. (2 $\Rightarrow$ 3): kekontinuan di $0$
dengan $\varepsilon = 1$ memberi $\delta$ dengan $\norm x \leq \delta
\Rightarrow \norm{u(x)} \leq 1$; lalu kehomogenannya menskalakan
sembarang $x$ dengan $\norm x \leq 1$ turun ke bola itu lalu kembali:
\[
\norm{u(x)} = \frac1\delta\,\norm{u(\delta x)} \leq
\frac1\delta ,
\]
sebab $\norm{\delta x} \leq \delta$. (3 $\Rightarrow$ 4):
untuk $x \neq 0$, terapkan batasnya pada $\frac{x}{\norm x}$. (4
$\Rightarrow$ 5): $\norm{u(x) - u(y)} = \norm{u(x - y)} \leq
C\norm{x - y}$. (5 $\Rightarrow$ 1) sudah diketahui.

Aksioma \omterm{def:b2:nvs:norm}{norma} bagi $\vertiii\cdot$: kehomogenan dan pemisahannya
jelas (sebab $\vertiii u = 0$ memaksa $u = 0$ pada bolanya, jadi di
mana-mana); ketaksamaan segitiganya dari $\norm{(u + v)(x)} \leq
\norm{u(x)} + \norm{v(x)}$. Sifat submultiplikatifnya: $\norm{v(u(x))}
\leq \vertiii v\,\norm{u(x)} \leq \vertiii v \vertiii u \norm x$.
\end{proof}

\begin{example}\label{ex:b2:nvs:operatornorms}
Pada $\bigl(C(\intcc{0}{1}), \norm\cdot_\infty\bigr)$: evaluasi
$f \mapsto f(0)$ bernorma operator $1$; pengintegralan $f \mapsto
\int_0^1 f$ bernorma $1$; sedangkan pemetaan $f \mapsto \int_0^1 t f(t)\dd t$
bernorma $\int_0^1 t\,\dd t = \frac12$ (batas atasnya dari ketaksamaan
segitiga untuk integral; dan tercapai di $f \equiv 1$). Namun
penurunan, dari $(C^1, \norm\cdot_\infty)$ ke $(C^0,
\norm\cdot_\infty)$, \emph{tidak} \omterm{def:b2:metric:continuity}{kontinu}: sebab $\norm{\sin(nx)}_\infty
= 1$ padahal turunannya bernorma supremum $n$. Linear tidak berarti
\omterm{def:b2:metric:continuity}{kontinu} dalam dimensi tak hingga.
\end{example}

\begin{example}[Dua norma, dua vonis atas satu barisan]\label{ex:b2:nvs:sqrtnxn}
Pada $C(\intcc01)$, misalkan $g_n(x) = \sqrt{n}\,x^n$. Maka
\[
\norm{g_n}_1 = \frac{\sqrt n}{n + 1} \longrightarrow 0,
\qquad
\norm{g_n}_2^2 = \frac{n}{2n + 1} \longrightarrow \frac12,
\qquad
\norm{g_n}_\infty = \sqrt n \longrightarrow \infty :
\]
satu barisan, tiga \omterm{def:b2:nvs:norm}{norma}, tiga perilaku --- konvergen ke
nol, tidak konvergen (\omterm{def:b2:nvs:norm}{normanya} mapan di $\frac1{\sqrt2}$ padahal
limit titik demi titiknya $0$), dan meledak. Massa yang memusat
di dekat $x = 1$ tak terlihat oleh $\norm\cdot_1$, setengah terlihat oleh
$\norm\cdot_2$, dan mendominasi bagi $\norm\cdot_\infty$. Dalam dimensi
tak hingga, ``apakah ia konvergen?'' bukan pertanyaan tentang sebuah
barisan: ia pertanyaan tentang sebuah barisan \emph{beserta} \omterm{def:b2:nvs:norm}{normanya}.
\end{example}

\begin{method}[Menghitung sebuah norma operator]\label{met:b2:nvs:opnorm}
Selalu dalam dua gerakan. \emph{Batas atas:} taksirlah
$\norm{u(x)}$ oleh $C\norm x$ dengan ketaksamaan segitiga,
Cauchy--Schwarz, atau batas integral --- ini membuktikan $\vertiii u
\leq C$. \emph{Saksi:} tunjukkan entah suatu $x_0 \neq 0$ tertentu
dengan $\norm{u(x_0)} = C\norm{x_0}$ (batasnya tercapai), atau
barisan vektor satuan $x_n$ dengan $\norm{u(x_n)} \to C$ (batasnya
dihampiri). Kedua gerakan itu wajib: batas atas semata hanya memberi
$\vertiii u \leq C$, sedangkan saksi semata hanya memberi
$\vertiii u \geq C$. Dalam dimensi tak hingga saksinya boleh jadi harus
berupa barisan --- sebab supremumnya tak harus tercapai
(\cref{exo:b2:nvs:8}).
\end{method}

\begin{example}[Operator diagonal melihat semua norma sama saja]\label{ex:b2:nvs:diagonalnorm}
Untuk $D = \operatorname{diag}(d_1, \dots, d_n)$ pada $K^n$ dengan
salah satu \emph{mana pun} dari \omterm{def:b2:nvs:norm}{norma} $\norm\cdot_1, \norm\cdot_2,
\norm\cdot_\infty$: dari $\abs{d_ix_i} \leq
\bigl(\max_j\abs{d_j}\bigr)\abs{x_i}$ koordinat demi koordinat diperoleh
$\norm{Dx} \leq \max_j\abs{d_j}\,\norm x$; dan $x = e_{j_0}$
(dengan indeks yang memaksimalkan) mencapainya. Jadi $\vertiii D =
\max_j\abs{d_j}$ pada ketiga kasusnya: bagi pemetaan diagonal, semua
\omterm{def:b2:nvs:norm}{norma} yang masuk akal menceritakan hal yang sama, yakni faktor
peregangan terbesar. Segala kesulitan tentang \omterm{thm:b2:nvs:continuouslinear}{norma operator} justru
berasal dari perilaku yang \emph{tidak} diagonal --- dan itulah sebabnya
\omterm{def:b2:nvs:norm}{norma} yang diselaraskan pada soal akhir pekan bab ini (pertanyaan 22)
bekerja dengan lebih dulu memaksa sebuah matriks menjadi diagonal.
\end{example}

\begin{example}[Jumlah kolom: kembaran \cref{exo:b2:nvs:4} untuk norma $1$]\label{ex:b2:nvs:columnsums}
Pada $(\R^n, \norm\cdot_1)$, \omterm{thm:b2:nvs:continuouslinear}{norma operator} sebuah matriks $A$ adalah
jumlah mutlak \emph{kolom} yang terbesar. Jalankan metodenya: untuk
$\norm x_1 \leq 1$,
\[
\norm{Ax}_1 = \sum_i\Bigl|\sum_j a_{ij}x_j\Bigr|
\leq \sum_j \abs{x_j}\sum_i\abs{a_{ij}}
\leq \Bigl(\max_j\sum_i\abs{a_{ij}}\Bigr)\norm x_1 ,
\]
dan batasnya tercapai di $x = e_{j_0}$ untuk kolom $j_0$ yang
memaksimalkan --- saksi paling rapi yang dapat dibayangkan. Jadi untuk $A =
\begin{psmallmatrix}1 & -2\\ 3 & 1\end{psmallmatrix}$:
$\vertiii A_1 = \max(1 + 3,\ 2 + 1) = 4$, sedangkan
$\vertiii A_\infty = 4$ juga (lewat barisnya) --- di sini kebetulan, bukan
hukum: ubahlah unsur matriksnya secara tidak setangkup dan kedua
\omterm{def:b2:nvs:norm}{normanya} akan berpisah. Baris untuk $\norm\cdot_\infty$, kolom untuk
$\norm\cdot_1$: jembatan keledainya, tiap \omterm{def:b2:nvs:norm}{norma} punya vektor satuan
(pola tanda, atau vektor basis) yang memungut jumlah yang bersesuaian.
\end{example}

\begin{proposition}[Pemetaan bilinear]\label{prop:b2:nvs:bilinear}
Pemetaan bilinear $b \colon E \times F \to G$ bersifat \omterm{def:b2:metric:continuity}{kontinu} bila dan
hanya bila $\norm{b(x,y)} \leq C\norm x\,\norm y$ untuk suatu $C$; dan
ketika itu ia \omterm{def:b2:metric:continuity}{Lipschitz} pada himpunan terbatas. (Pola buktinya sama;
hasil kali $(u, v) \mapsto v \circ u$ dan perkalian matriks adalah contoh
kuncinya.)
\end{proposition}

\begin{proof}
Jika batas itu berlaku, maka
\[
b(x,y) - b(x_0,y_0) = b(x - x_0,\, y) + b(x_0,\, y - y_0),
\]
sehingga $\norm{b(x,y) - b(x_0,y_0)} \leq C\norm{x - x_0}\norm y +
C\norm{x_0}\norm{y - y_0}$: itulah kekontinuan di $(x_0, y_0)$, beserta
batas \omterm{def:b2:metric:continuity}{Lipschitz} di tempat $\norm x, \norm y \leq R$. Sebaliknya,
kekontinuan di $(0,0)$ memberi $\delta$ dengan $\norm{b(x,y)} \leq 1$ pada
$\norm x, \norm y \leq \delta$; lalu skalakan kedua peubahnya.
\end{proof}

\section{Dimensi hingga}

\begin{theorem}[Kesetaraan norma dalam dimensi hingga]\label{thm:b2:nvs:finitedim}
Pada ruang berdimensi hingga, \emph{semua \omterm{def:b2:nvs:equivalent}{norma
setara}}.\index{kesetaraan norma} Akibatnya, dalam dimensi
hingga: kekonvergenan, keterbukaan, kekompakan, dan kelengkapan menjadi
gagasan yang tak bergantung pada \omterm{def:b2:nvs:norm}{normanya}; \omterm{def:b2:metric:compact}{kompak} $=$ tertutup dan
terbatas; ruangnya \omterm{def:b2:metric:complete}{lengkap}; dan setiap pemetaan linear (atau
multilinear) \emph{dari} ruang berdimensi hingga bersifat \omterm{def:b2:metric:continuity}{kontinu}.
\end{theorem}

\begin{proof}
Tetapkan sebuah basis lalu samakan $E \simeq K^n$; cukuplah
membandingkan sembarang \omterm{def:b2:nvs:norm}{norma} $N$ dengan $\norm\cdot_\infty$.

\emph{Satu arahnya aljabar:} $N(x) = N(\sum x_i e_i) \leq
\sum \abs{x_i} N(e_i) \leq C \norm x_\infty$ dengan $C = \sum N(e_i)$.
Ini sekaligus menunjukkan $N$ \omterm{def:b2:metric:continuity}{kontinu} pada $(K^n, \norm\cdot_\infty)$
(sebab ia Lipschitz-$C$: $\abs{N(x) - N(y)} \leq N(x - y)$).

\emph{Arah satunya topologi:} permukaan bola satuan $S = \{x :
\norm{x}_\infty = 1\}$ bersifat tertutup dan terbatas di $(K^n,
\norm\cdot_\infty)$, jadi \omterm{def:b2:metric:compact}{kompak}
(\cref{thm:b2:metric:compactprops}~(2), yang berlaku untuk $\C^n \simeq
\R^{2n}$). Fungsi \omterm{def:b2:metric:continuity}{kontinu} $N$ mencapai minimumnya $c$ pada
$S$; dan $c > 0$ sebab $N$ hanya nol di $0 \notin S$. Kehomogenannya
lalu menyebarkan batas itu: $N(x) \geq c \norm{x}_\infty$ untuk setiap $x$.

Akibatnya: semua pernyataannya menyusut ke $(K^n,
\norm\cdot_\infty)$, tempat semuanya sudah diketahui
(\cref{thm:b2:metric:rncomplete},
\cref{thm:b2:metric:compactprops}); dan pemetaan linear $u$ dari
$E$ yang berdimensi hingga memenuhi $\norm{u(x)} \leq \sum\abs{x_i}
\norm{u(e_i)} \leq C'\norm{x}_\infty$: yakni batas (4) pada
\cref{thm:b2:nvs:continuouslinear}.
\end{proof}

\begin{corollary}\label{cor:b2:nvs:closedsubspace}
Subruang berdimensi hingga pada ruang bernorma apa pun bersifat tertutup.
\end{corollary}

\begin{proof}
Ia \omterm{def:b2:metric:complete}{lengkap} terhadap \omterm{def:b2:nvs:norm}{norma} imbasannya
(\cref{thm:b2:nvs:finitedim}), sedangkan himpunan bagian yang \omterm{def:b2:metric:complete}{lengkap}
bersifat tertutup (\cref{def:b2:metric:complete}).
\end{proof}

\begin{example}[Sebuah hampiran terbaik yang dihitung lewat kesetangkupan]\label{ex:b2:nvs:bestaffine}
Di $\bigl(C(\intcc{-1}{1}), \norm\cdot_\infty\bigr)$, seberapa jauh
$f(x) = \abs x$ dari subruang berdimensi dua yang tertutup berisi
fungsi afin $a + bx$? Berkat kesetangkupannya, mengganti $a + bx$ dengan
$a - bx$ tidak mengubah $\norm{f - (a \pm bx)}_\infty$, dan
titik tengahnya $a$ setidaknya sama baiknya (lewat ketaksamaan segitiga
pada rata-ratanya): jadi cukuplah meninjau fungsi konstan. Untuk
konstanta $a$:
\[
\norm{\abs x - a}_\infty = \max\,(1 - a,\ a)
\geq \frac12 ,
\]
yang minimum di $a = \frac12$: jadi jaraknya $\frac12$, tercapai
oleh konstanta $\frac12$. Perhatikan kurva galatnya $\abs x -
\frac12$: ia mencapai $\pm\frac12$ berselang-seling di $x = -1, 0,
1$ --- tiga ekstremum bertanda berselang-seling bagi hampiran
terbaik dari keluarga berparameter dua. Pola
\emph{ekuiosilasi} itu bukan kebetulan; ia sidik jari keoptimalan yang
diubah soal akhir pekan bab ini menjadi teorema Chebyshev.
\end{example}

\begin{example}[Subruang tertutup lawan subruang padat]\label{ex:b2:nvs:closeddense}
Di $E = \bigl(C(\intcc{0}{1}), \norm\cdot_\infty\bigr)$: tiap
$\R_n[X]$ (polinomial berderajat $\leq n$, dibatasi pada
$\intcc01$) merupakan subruang berdimensi hingga, jadi \emph{tertutup}
--- sebab limit seragam polinomial berderajat $\leq n$ tetap demikian.
Namun gabungan $\R[X]$ atas semuanya bersifat \emph{padat} di $E$
(lewat teorema hampiran Weierstrass, yang dibuktikan pada
\cref{ch:b2:funcseq}), dan subruang sejati yang padat sama sekali tidak
tertutup. Pelajarannya: ketertutupan subruang adalah keistimewaan
dimensi hingga; menumpuk lantai yang tertutup dapat membangun
pencakar langit yang padat.
\end{example}

\begin{theorem}[Riesz]\label{thm:b2:nvs:riesz}
Bola satuan tertutup sebuah ruang bernorma $E$ bersifat \omterm{def:b2:metric:compact}{kompak}
\emph{bila dan hanya bila} $\dim E < \infty$.\index{teorema Riesz}
\end{theorem}

\begin{proof}
Untuk dimensi hingga: cukuplah tertutup dan terbatas
(\cref{thm:b2:nvs:finitedim}).

Sebaliknya, andaikan $\dim E = \infty$. \emph{Lema Riesz:} untuk
setiap subruang tertutup sejati $F \subsetneq E$ dan $\varepsilon \in
\intoo{0}{1}$, ada vektor satuan $x$ dengan $d(x, F) \geq 1 -
\varepsilon$. Buktinya: pilih $y \notin F$, misalkan $\delta = d(y, F) > 0$
(sebab $F$ tertutup), pilihlah $f \in F$ dengan $\norm{y - f} \leq
\frac{\delta}{1 - \varepsilon}$, lalu tetapkan $x = \frac{y - f}{\norm{y -
f}}$: maka untuk sembarang $g \in F$,
\[
\norm{x - g} = \frac{\norm{y - (f + \norm{y-f}\,g)}}{\norm{y - f}}
\geq \frac{\delta}{\norm{y-f}} \geq 1 - \varepsilon ,
\]
sebab pembilangnya adalah jarak dari $y$ ke suatu titik $F$.

Kini bangunlah vektor satuan $x_1, x_2, \dots$ secara induktif: $F_k =
\operatorname{Vect}(x_1, \dots, x_k)$ berdimensi hingga, jadi
tertutup (\cref{cor:b2:nvs:closedsubspace}) dan sejati; lalu lema Riesz
dengan $\varepsilon = \frac12$ menyediakan vektor satuan $x_{k+1}$ dengan
$d(x_{k+1}, F_k) \geq \frac12$. Barisan itu memenuhi $\norm{x_p - x_q}
\geq \frac12$ untuk $p \neq q$: jadi tak ada barisan bagian yang
konvergen --- sehingga bola satuannya tidak \omterm{def:b2:metric:compact}{kompak}.
\end{proof}

\begin{example}[Riesz sebagai pendeteksi dimensi]\label{ex:b2:nvs:rieszdetector}
Apakah $C(\intcc01)$ berdimensi hingga? Riesz menjawabnya tanpa
menunjukkan satu pun keluarga bebas tak hingga yang gamblang: barisan
$f_n(x) = x^n$ berada di bola satuan tertutup dan memenuhi, untuk
$m > n$, $\norm{f_n - f_m}_\infty \geq f_n(x_0) - f_m(x_0) > 0$
di titik yang sesuai --- dan hal itu dikuantifikasi dengan rapi pada soal
akhir pekan bab ini (pertanyaan 16), tempat sebuah barisan bagian tetap
berjarak antara $\geq \frac14$. Jadi tak ada barisan bagian yang
konvergen, sehingga bolanya tidak \omterm{def:b2:metric:compact}{kompak}, sehingga $\dim C(\intcc01) =
\infty$ menurut \cref{thm:b2:nvs:riesz}. Kekompakan bola satuan adalah
dikotomi yang sempurna: ia berlaku dalam dimensi hingga, gagal dalam
dimensi tak hingga, tanpa jalan tengah --- geometrinya sendiri sudah
membaca jenis dimensinya.
\end{example}

\section{Ruang Banach}

\begin{definition}\label{def:b2:nvs:banach}
Sebuah \emph{ruang Banach}\index{ruang Banach} adalah ruang bernorma
yang \omterm{def:b2:metric:complete}{lengkap}. Contohnya: setiap ruang bernorma berdimensi hingga
(\cref{thm:b2:nvs:finitedim}); $\bigl(C(\intcc{a}{b}),
\norm\cdot_\infty\bigr)$ (\cref{thm:b2:metric:rncomplete});
$\mathcal{L}_c(E, F)$ untuk $F$ yang Banach (dengan pola bukti yang sama
seperti untuk fungsi \omterm{def:b2:metric:continuity}{kontinu}). Bukan contohnya: $\bigl(C(\intcc{0}{1}),
\norm\cdot_1\bigr)$ (\cref{exo:b2:nvs:7}).
\end{definition}

\begin{example}[Norma operator pengintegralan]\label{ex:b2:nvs:integrationnorm}
Pada $\bigl(C(\intcc01), \norm\cdot_\infty\bigr)$, misalkan $T(f)(x) =
\int_0^x f(t)\,\dd t$ (sebuah endomorfisma, sebab $T(f)$ \omterm{def:b2:metric:continuity}{kontinu}).
Jalankan \cref{met:b2:nvs:opnorm}. Batas atasnya:
\[
\abs{T(f)(x)} \leq \int_0^x\abs f \leq x\,\norm f_\infty \leq
\norm f_\infty ,
\]
jadi $\vertiii T \leq 1$. Saksinya: $f \equiv 1$ memberi $T(f)(x) =
x$ dan $\norm{T(f)}_\infty = 1 = \norm f_\infty$: jadi tercapai, dan
$\vertiii T = 1$. Namun perhatikan $\vertiii{T^2} = \frac12 \neq
\vertiii T^2$: memang $T^2(f)(x) = \int_0^x(x - t)f(t)\dd t$ punya
$\abs{T^2(f)(x)} \leq \frac{x^2}2\norm f_\infty$, yang lagi-lagi tercapai
di $f \equiv 1$; dan secara umum $\vertiii{T^n} = \frac1{n!}$ ---
sedangkan batas submultiplikatifnya $\vertiii T^n = 1$ meleset sejauh
satu faktorial. Justru gejala inilah yang diubah siasat iterat pada
soal akhir pekan \cref{ch:b2:metric} menjadi keterselesaian global
persamaan diferensial linear.
\end{example}

\begin{theorem}[Kekonvergenan mutlak pada ruang Banach]\label{thm:b2:nvs:absoluteconvergence}
Pada \omterm{def:b2:nvs:banach}{ruang Banach}, jika $\sum \norm{u_n} < \infty$ maka $\sum u_n$
konvergen, dan $\norm{\sum u_n} \leq \sum\norm{u_n}$. (Teori \omterm{def:b2:metric:complete}{lengkap}
deret pada ruang bernorma menjadi isi \cref{ch:b2:series}.)
\end{theorem}

\begin{proof}
Untuk jumlah parsial $S_N$: bila $q > p$, maka $\norm{S_q - S_p} \leq
\sum_{n=p+1}^{q}\norm{u_n}$, yang menuju $0$ (lewat kriteria Cauchy
untuk deret real berisi \omterm{def:b2:nvs:norm}{normanya}): jadi $(S_N)$ bersifat Cauchy,
sehingga konvergen. Ketaksamaannya berpindah ke limit dari ketaksamaan
segitiga yang hingga.
\end{proof}

\begin{example}[Eksponensial matriks, perkenalan pertama]\label{ex:b2:nvs:matrixexp}
\emph{Eksponensial matriks}\index{eksponensial matriks}:
$\mathcal{M}_n(K)$ dengan \omterm{def:b2:nvs:norm}{norma} submultiplikatif apa pun ($\vertiii{AB}
\leq \vertiii A \vertiii B$) bersifat Banach (karena berdimensi hingga).
Maka, untuk setiap $A$,
\[
\eu^A = \sum_{k=0}^{\infty} \frac{A^k}{k!}
\]
konvergen mutlak (sebab $\vertiii{A^k/k!} \leq \vertiii A^k /k!$, yang
terjumlahkan): jadi ia terdefinisi dengan baik. \cref{ch:b2:diffeq}
memanfaatkannya secara sistematis.
\end{example}

\begin{example}[Deret Neumann yang berhenti]\label{ex:b2:nvs:neumannfinite}
Untuk $A = \begin{psmallmatrix}0 & \frac12\\ 0 &
0\end{psmallmatrix}$: berlaku $\vertiii A < 1$ pada \omterm{thm:b2:nvs:continuouslinear}{norma operator} apa
pun yang \omterm{def:b2:structures:generated}{dibangun} atas \omterm{def:b2:nvs:norm}{norma} \cref{ex:b2:nvs:examples}, dan $A^2 = 0$,
sehingga deret geometrinya runtuh:
\[
(I - A)^{-1} = \sum_{k \geq 0} A^k = I + A =
\begin{pmatrix}1 & \tfrac12\\ 0 & 1\end{pmatrix},
\]
yang terperiksa lewat $(I - A)(I + A) = I - A^2 = I$. Kenilpotenan
memangkas deretnya persis seperti ia memangkas eksponensial pada
\cref{ch:b2:reduction}; dan contoh ini menakar
harapan kita: invers Neumann pada umumnya berupa deret tak hingga,
dan berupa polinomial persis ketika usikannya nilpoten,
sedangkan galat setelah $N$ suku selalu terbatas oleh
ekor geometrinya $\vertiii A^{N+1}/(1 - \vertiii A)$.
\end{example}

\begin{example}[Eksponensial sebuah pembangkit rotasi]\label{ex:b2:nvs:rotationexp}
Ambil $A = \begin{psmallmatrix}0 & -\theta\\ \theta &
0\end{psmallmatrix}$. Maka $A^2 = -\theta^2 I$, sehingga pangkatnya
berputar dengan periode empat, dan deretnya terpecah menjadi bagian genap
dan bagian ganjil:
\[
\eu^{A} = \sum_{k}\frac{A^k}{k!}
= \Bigl(\sum_{j}\frac{(-1)^j\theta^{2j}}{(2j)!}\Bigr) I
+ \Bigl(\sum_{j}\frac{(-1)^j\theta^{2j+1}}{(2j+1)!}\Bigr)
\frac{A}{\theta}
= \begin{pmatrix}
\cos\theta & -\sin\theta\\
\sin\theta & \cos\theta
\end{pmatrix},
\]
dengan semua penataan ulangnya disahkan oleh kekonvergenan mutlak.
Eksponensial pembangkit yang antisetangkup adalah sebuah rotasi ---
yang di sini dihitung murni dari deretnya, tiga bab sebelum
persamaan diferensial $x' = Ax$ (\cref{ch:b2:diffeq}) menjelaskan
\emph{mengapa}: sebab $\eu^{tA}$ adalah gerak melingkar beraturan.
Pelajaran penutupnya: kesamaan antar deret matriks dibuktikan persis
seperti kesamaan skalar, asalkan sebuah \omterm{def:b2:nvs:norm}{norma} submultiplikatif
mengesahkan kekonvergenan mutlaknya.
\end{example}

\begin{example}[Norma supremum berarti seragam: kamusnya]\label{ex:b2:nvs:uniformdictionary}
Pernyataan $\norm{f_n - f}_\infty \to 0$ \emph{adalah} kekonvergenan
seragam: satu bilangan, yakni $\sup_x\abs{f_n(x) - f(x)}$, membatasi
galatnya di setiap titik sekaligus. Kamus itu bekerja pada $f_n(x) = x^n$
atas $\intcc{0}{1}$: titik demi titik, $f_n \to 0$ pada $\intco{0}{1}$
sedangkan $f_n(1) = 1$; dalam \omterm{def:b2:nvs:norm}{norma}, $\norm{f_n - 0}_\infty = 1 \not\to
0$, dan memang limit titik demi titiknya tidak \omterm{def:b2:metric:continuity}{kontinu}, jadi mustahil
menjadi limit terhadap $\norm\cdot_\infty$ di $C(\intcc01)$ (yang
tertutup terhadap limit seragam, \cref{thm:b2:metric:rncomplete}). Pada
$\intcc{0}{a}$ dengan $a < 1$: $\norm{f_n}_\infty = a^n \to 0$ ---
kekonvergenan seragamnya pulih setelah daerahnya diperkecil. Setiap
pernyataan kekonvergenan pada \cref{ch:b2:funcseq} adalah pernyataan
tentang satu \omterm{def:b2:nvs:norm}{norma} ini; mengingat kamusnya memaruhkan bab tersebut.
\end{example}

\begin{remark}[Jebakan yang sering muncul]\label{rem:b2:nvs:pitfalls}
(i) \omterm{thm:b2:nvs:continuouslinear}{Norma operator} bergantung pada \emph{kedua} \omterm{def:b2:nvs:norm}{norma} yang dipilih:
matriks yang sama punya $\vertiii\cdot_\infty$ berupa jumlah baris
(\cref{exo:b2:nvs:4}) dan $\vertiii\cdot_1$ yang berbeda (jumlah
kolom); jadi menyebut ``\omterm{def:b2:nvs:norm}{norma}'' sebuah matriks tanpa menamai
\omterm{def:b2:nvs:norm}{norma} yang melandasinya tidak bermakna apa pun. (ii) $\vertiii{AB} \leq
\vertiii A\,\vertiii B$ adalah ketaksamaan, dan lazimnya tegas ---
pangkatnya dapat menyusut jauh lebih cepat daripada dugaan batas
$\vertiii A^k$, dan justru itulah inti \omterm{def:b2:nvs:norm}{norma} yang diselaraskan (soal
akhir pekan bab ini, pertanyaan 22). (iii) ``Linear berarti
\omterm{def:b2:metric:continuity}{kontinu}'' adalah keistimewaan dimensi hingga: penurunan
pada polinomial bersifat linear dan tak terbatas
(\cref{ex:b2:nvs:operatornorms}). (iv) Kekonvergenan mutlak
$\sum u_n$ hanya menolong bila ruangnya \omterm{def:b2:metric:complete}{lengkap}
(\cref{exo:b2:series:9} membangun contoh penyangkalnya). (v) Dalam
dimensi tak hingga, supremum atas bola satuan itu sungguh-sungguh
supremum: jangan mengandaikan ia tercapai
(\cref{exo:b2:nvs:8}).
\end{remark}

\begin{remark}[Pandangan ke depan di dalam jilid ini]\label{rem:b2:nvs:perspectives}
Kini tiga janji temu sudah tertetapkan. Dengan \cref{ch:b2:series}: pada
\omterm{def:b2:nvs:banach}{ruang Banach}, deret yang konvergen mutlak pasti konvergen, sehingga
deret geometri dan deret eksponensial atas operator menjadi
perkakas sehari-hari --- untuk membalik $I - A$ dan menetapkan $\eu^{A}$
(\cref{ex:b2:series:neumann}). Dengan
\cref{ch:b2:funcseq} dan \cref{ch:b2:powerseries}: kekonvergenan
barisan fungsi dan deret pangkat tak lain kekonvergenan di
$\bigl(C, \norm\cdot_\infty\bigr)$
(\cref{ex:b2:nvs:uniformdictionary}), dan jari-jari
kekonvergenannya adalah pernyataan tentang deret geometri mana yang
mendominasi. Dengan \cref{ch:b2:fourier}: \omterm{def:b2:nvs:norm}{norma}
$\norm\cdot_2$ dan $\norm\cdot_\infty$ sungguh berselisih pada
$C(\intcc{0}{1})$ (\cref{ex:b2:nvs:sqrtnxn}), dan persis itulah
sebabnya kekonvergenan rata-rata kuadrat deret Fourier dan kekonvergenan
seragam menjadi dua teorema berbeda dengan harga yang berbeda pula.
\end{remark}

\begin{remark}[Di mana bab ini dipakai]
\omterm{thm:b2:nvs:continuouslinear}{Norma operator} dan deret geometri menggerakkan hujah usikan
pada \cref{ch:b2:diffcalc} (teorema fungsi invers) dan
\omterm{ex:b2:nvs:matrixexp}{eksponensial matriks} pada \cref{ch:b2:diffeq}; \omterm{thm:b2:nvs:finitedim}{kesetaraan
norma} diam-diam mengesahkan setiap hujah ``pilih saja \omterm{def:b2:nvs:norm}{norma}
kesukaanmu'' pada \cref{ch:b2:funcseq} dan seterusnya; sedangkan
jurang hingga/tak hingga pada teorema Riesz --- yang dikuantifikasi pada
soal akhir pekan bab ini --- adalah alasan mengapa jilid Tahun
ke-3 harus memakai perkakas baru (kekonvergenan lemah, Arzelà--Ascoli,
proyeksi pada ruang Hilbert) di tempat jilid ini masih dapat menyarikan
barisan bagian yang konvergen.
\end{remark}

\section{Latihan}

\begin{exercise}[$\star$]\label{exo:b2:nvs:1}
Pada $\R^2$, gambarlah bola satuan $\norm\cdot_1$, $\norm\cdot_2$,
$\norm\cdot_\infty$, lalu buktikan ketaksamaan $\norm x_\infty
\leq \norm x_2 \leq \norm x_1 \leq 2\norm x_\infty$ dengan konstanta
terbaik pada dimensi $2$.
\end{exercise}

\begin{exercise}[$\star$]\label{exo:b2:nvs:2}
Apakah $N(f) = \abs{f(0)} + \norm{f'}_\infty$ merupakan \omterm{def:b2:nvs:norm}{norma} pada
$C^1(\intcc{0}{1})$? Bandingkan dengan $\norm{f}_\infty$: satu
ketaksamaannya berlaku, satunya gagal (tunjukkan).
\end{exercise}

\begin{exercise}[$\star$]\label{exo:b2:nvs:3}
Hitunglah \omterm{thm:b2:nvs:continuouslinear}{norma operator} $u(f) = \int_0^1 f(t)\,\eu^t\,\dd t$
pada $\bigl(C(\intcc{0}{1}), \norm\cdot_\infty\bigr) \to \R$, dan \omterm{thm:b2:nvs:continuouslinear}{norma
operator} geseran $S(x_1, x_2, \dots, x_n) = (x_2, \dots, x_n, 0)$ pada
$(K^n, \norm\cdot_\infty)$.
\end{exercise}

\begin{exercise}[$\star\star$]\label{exo:b2:nvs:4}
Pada $(\R^n, \norm\cdot_\infty)$, buktikan bahwa \omterm{thm:b2:nvs:continuouslinear}{norma operator} sebuah
matriks $A$ adalah $\vertiii A_\infty = \max_i \sum_j \abs{a_{ij}}$
(yakni jumlah mutlak baris yang terbesar). Hitunglah untuk $\begin{pmatrix} 1 & -2\\
3 & 1\end{pmatrix}$.
\end{exercise}

\begin{exercise}[$\star\star$]\label{exo:b2:nvs:5}
Buktikan bahwa $GL_n(K)$ \omterm{def:b2:metric:topology}{terbuka} di $\mathcal{M}_n(K)$ dan bahwa $A
\mapsto A^{-1}$ \omterm{def:b2:metric:continuity}{kontinu} di sana. \emph{Petunjuk: untuk keterbukaannya,
jika $\vertiii H < \frac{1}{\vertiii{A^{-1}}}$ maka $A + H = A(I +
A^{-1}H)$ dengan $\vertiii{A^{-1}H} < 1$, sedangkan $I + B$ punya invers
untuk $\vertiii B < 1$ berkat deret geometrinya
(\cref{thm:b2:nvs:absoluteconvergence}); untuk kekontinuannya, batasi
$(A+H)^{-1} - A^{-1}$ dengan deret yang sama.}
\end{exercise}

\begin{exercise}[$\star\star$]\label{exo:b2:nvs:6}
Misalkan $\varphi$ sebuah bentuk linear pada ruang bernorma $E$.
Buktikan bahwa $\varphi$ \omterm{def:b2:metric:continuity}{kontinu} bila dan hanya bila $\ker\varphi$
tertutup. \emph{(Jika $\ker\varphi$ tertutup dan $\varphi \neq 0$, pilih
$a$ dengan $\varphi(a) = 1$ dan $r > 0$ dengan $B(a, r) \cap \ker\varphi
= \emptyset$; lalu turunkan $\abs{\varphi(h)} \leq \frac{1}{r}\norm h$
lewat hujah penskalaan pada $a - \frac{h}{\varphi(h)}$.)}
\end{exercise}

\begin{exercise}[$\star\star$]\label{exo:b2:nvs:7}
Buktikan bahwa $\bigl(C(\intcc{0}{1}), \norm\cdot_1\bigr)$ tidak
\omterm{def:b2:metric:complete}{lengkap}: tunjukkan bahwa fungsi $f_n$, berupa tanjakan afin dari $0$ ke
$1$ pada $\bigl[\frac12 - \frac1n, \frac12\bigr]$ (bernilai $0$
sebelumnya dan $1$ sesudahnya), membentuk barisan Cauchy yang tak punya
limit \omterm{def:b2:metric:continuity}{kontinu} terhadap $\norm\cdot_1$.
\end{exercise}

\begin{exercise}[$\star\star\star$]\label{exo:b2:nvs:8}
Pada $E = C(\intcc{0}{1})$ dengan $\norm\cdot_\infty$, tinjaulah
\[
\varphi(f) = \sum_{n \geq 1} (-1)^n\, 2^{-n} f\bigl(\tfrac1n\bigr).
\]
Buktikan bahwa $\varphi$ merupakan bentuk linear \omterm{def:b2:metric:continuity}{kontinu} yang
terdefinisi dengan baik dengan $\vertiii\varphi = 1$, tetapi supremum
yang mendefinisikan $\vertiii\varphi$ \emph{tidak tercapai} pada bola
satuan tertutupnya. \emph{(Batas atasnya: ketaksamaan segitiga. \omterm{def:b2:nvs:norm}{Normanya}
$= 1$: bangunlah $f_K$ yang \omterm{def:b2:metric:continuity}{kontinu} dengan $\norm{f_K}_\infty \leq 1$ dan
$f_K(\frac1n) = (-1)^n$ untuk $n \leq K$ --- sebab titik $\frac1n$
terpencil satu sama lain. Ketaktercapaiannya: kesamaan akan memaksa
$f(\frac1n) = (-1)^n$ untuk setiap $n$, dan itu tak selaras dengan
kekontinuan $f$ di $0$ karena $\frac1n \to 0$.)}
\end{exercise}

\begin{exercise}[$\star\star\star$]\label{exo:b2:nvs:9}
Misalkan $E$ ruang bernorma yang bola satuan tertutupnya \omterm{def:b2:metric:compact}{kompak}.
Turunkan kembali, tanpa mengutip \cref{thm:b2:nvs:riesz}, bahwa setiap
barisan terbatas punya barisan bagian yang konvergen, lalu buktikan bahwa
setiap bentuk linear pada $E$ \omterm{def:b2:metric:continuity}{kontinu} bila dan hanya bila $\dim E <
\infty$. \emph{(Untuk dimensi tak hingga, bangunlah bentuk yang tak
\omterm{def:b2:metric:continuity}{kontinu} dengan menetapkannya secara bebas pada barisan bebas linear yang
ternormalkan lalu memperluasnya --- dengan menerima keberadaan pelengkap
aljabarnya.)}
\end{exercise}

\begin{exercise}[$\star\star$]\label{exo:b2:nvs:10}
Pada $C(\intcc{0}{1})$, buktikan $\norm f_1 \leq \norm f_2 \leq \norm
f_\infty$ \emph{(lewat Cauchy--Schwarz untuk yang pertama)}, lalu
tunjukkan dengan keluarga $f_n(x) = x^n$ bahwa tak satu pun
ketaksamaannya dapat dibalik hingga sebuah konstanta: jadi ketiga
\omterm{def:b2:nvs:norm}{normanya} tak \omterm{def:b2:nvs:equivalent}{setara} sepasang demi sepasang.
\end{exercise}

\begin{exercise}[$\star\star$]\label{exo:b2:nvs:11}
(Jarak ke sebuah hiperbidang) Misalkan $\varphi$ bentuk linear \omterm{def:b2:metric:continuity}{kontinu}
yang tak nol pada ruang bernorma $E$. Buktikan bahwa
\[
d\bigl(x, \ker\varphi\bigr) =
\frac{\abs{\varphi(x)}}{\vertiii\varphi}
\qquad (x \in E),
\]
lalu periksa pada \cref{exo:b2:nvs:8} bahwa infimumnya tak harus
tercapai oleh titik hiperbidangnya.
\end{exercise}

\begin{exercise}[$\star\star\star$]\label{exo:b2:nvs:12}
Pada $E = \R[X]$ (semua polinomial), misalkan $N_1(P) =
\sup_{\intcc{0}{1}}\abs P$ dan $N_2(P) =
\sup_{\intcc{0}{2}}\abs P$. Tunjukkan bahwa $N_1 \leq N_2$ tetapi
$N_1$ dan $N_2$ \emph{tidak} \omterm{def:b2:nvs:equivalent}{setara}; lalu turunkan bahwa
identitas $(E, N_2) \to (E, N_1)$ merupakan bijeksi linear \omterm{def:b2:metric:continuity}{kontinu}
yang inversnya tidak \omterm{def:b2:metric:continuity}{kontinu}. Terakhir tunjukkan bahwa $(E, N_1)$
tidak \omterm{def:b2:metric:complete}{lengkap} \emph{(lewat jumlah parsial Taylor $\eu^x$)}. Ketiga
gejala itu mustahil dalam dimensi hingga --- katakanlah mengapa.
\end{exercise}

\section{Soal: Hampiran Terbaik dan Teorema Chebyshev}

Seberapa baik sebuah fungsi dapat dihampiri oleh polinomial berderajat
tertentu, dan polinomial mana yang paling baik? Untuk sisi keberadaannya,
jawabannya milik bab ini: kekompakan dalam dimensi
hingga membuat hampiran terbaik itu ada. Untuk sisi gamblangnya,
satu kasus tak sepele dapat diselesaikan tuntas dengan tangan kosong ---
di antara semua polinomial \emph{monik} berderajat $n$, yang bernorma
supremum terkecil pada $\intcc{-1}{1}$ adalah polinomial Chebyshev
(yang dinormalkan), dengan \omterm{def:b2:nvs:norm}{norma} $2^{1-n}$: itulah \emph{teorema
ekstremal Chebyshev}. Soal ini membuktikan kedua sisinya, lalu menakar
seberapa parah kekompakan itu gagal dalam dimensi tak hingga: bola satuan
$C(\intcc{0}{1})$ memuat gugusan titik tak hingga banyak yang jarak
antaranya $1$.

\begin{problem}[{Soal akhir pekan --- teorema ekstremal Chebyshev
dan geometri bola satuan}]
\label{pb:b2:nvs:1}
\omterm{def:b2:nvs:norm}{Norma} tanpa indeks adalah \omterm{def:b2:nvs:norm}{norma} supremum pada ruas yang disebutkan.

\medskip
\noindent\textbf{Bagian I --- Hampiran terbaik pada ruang
bernorma.}
\begin{enumerate}
  \item Misalkan $F$ subruang berdimensi hingga sebuah ruang
        bernorma $E$ dan $x \in E$. Buktikan bahwa jarak $d(x,
        F) = \inf_{f \in F}\norm{x - f}$ \emph{tercapai}
        \emph{(susutkan ke himpunan bagian $F$ yang tertutup dan terbatas
        lalu pakai \cref{thm:b2:nvs:finitedim})}.
  \item Sebuah \omterm{def:b2:nvs:norm}{norma} disebut \emph{cembung tegas} apabila $\norm u =
        \norm v = 1$ dan $u \neq v$ mengakibatkan $\bigl\Vert\frac{u +
        v}2\bigr\Vert < 1$. Tunjukkan bahwa $\norm\cdot_2$ pada $\R^n$
        bersifat cembung tegas \emph{(lewat kesamaan jajargenjang)}, dan
        bahwa $\norm\cdot_1$ dan $\norm\cdot_\infty$ tidak demikian untuk
        $n \geq 2$.
  \item Buktikan bahwa untuk \omterm{def:b2:nvs:norm}{norma} yang cembung tegas, hampiran
        terbaik pada pertanyaan 1 bersifat \emph{tunggal}.
  \item Di $(\R^2, \norm\cdot_\infty)$, hitunglah semua hampiran
        terbaik bagi $x = (0, 1)$ oleh garis $F =
        \operatorname{Vect}\bigl((1,0)\bigr)$: yakni satu interval
        berisi pemberi minimum.
  \item Di $\bigl(C(\intcc{a}{b}), \norm\cdot_\infty\bigr)$,
        tunjukkan bahwa hampiran terbaik $f$ oleh fungsi
        \emph{konstan} bersifat tunggal, sama dengan $c^* = \frac{\max f
        + \min f}{2}$, dengan jarak $\frac{\max f - \min
        f}{2}$; lalu hitunglah keduanya untuk $f(x) = x^2$ pada
        $\intcc{0}{1}$.
\end{enumerate}

\medskip
\noindent\textbf{Bagian II --- Polinomial Chebyshev.}
\begin{enumerate}[resume]
  \item Tunjukkan bahwa ada tepat satu polinomial $T_n$ dengan
        $T_n(\cos\theta) = \cos n\theta$ untuk setiap $\theta$
        \emph{(lewat rekurensi $T_{n+1} = 2XT_n - T_{n-1}$ dari
        rumus penjumlahan kosinus)}, bahwa $\deg T_n = n$, dan bahwa
        koefisien utamanya $2^{n-1}$ untuk $n \geq 1$.
  \item Tunjukkan $\abs{T_n} \leq 1$ pada $\intcc{-1}{1}$, dengan
        $T_n(\eta_k) = (-1)^k$ pada $n + 1$ titik $\eta_k =
        \cos\frac{k\pi}{n}$ ($k = 0, \dots, n$), dan bahwa
        akar $T_n$ adalah $n$ titik
        $\cos\frac{(2k-1)\pi}{2n}$, yang menyelang-nyeling $\eta_k$.
  \item Hitunglah $T_2, T_3, T_4$, lalu periksalah pergantian tanda
        $T_3$ di $\eta_0, \dots, \eta_3 = 1, \frac12, -\frac12,
        -1$ lewat penilaian langsung.
  \item Untuk $\abs x \geq 1$, buktikan
        \[
        T_n(x) = \frac{\bigl(x + \sqrt{x^2 - 1}\bigr)^n +
        \bigl(x - \sqrt{x^2 - 1}\bigr)^n}{2},
        \]
        lalu turunkan $T_n(x) \sim \frac12\bigl(x + \sqrt{x^2 -
        1}\bigr)^n \to \infty$ secara geometri untuk $x > 1$ yang tetap.
  \item Buktikan hukum penyusunan $T_m \circ T_n = T_{mn}$
        \emph{(periksa pada $\intcc{-1}{1}$ lalu panggil kekakuan
        polinomial)}.
\end{enumerate}

\medskip
\noindent\textbf{Bagian III --- Teorema ekstremal Chebyshev.}
Tulis $Q_n = 2^{1-n}T_n$ (yang monik, menurut pertanyaan 6).
\begin{enumerate}[resume]
  \item Misalkan $P$ monik berderajat $n \geq 1$ dengan
        $\sup_{\intcc{-1}{1}}\abs P < 2^{1-n}$. Dengan menilai $D
        = Q_n - P$ di titik $\eta_k$ lalu mencacah pergantian
        tandanya, turunkan sebuah kontradiksi. Simpulkan:
        \[
        \sup_{\intcc{-1}{1}}\abs P \;\geq\; 2^{1-n}
        \qquad\text{untuk setiap } P \text{ monik berderajat } n.
        \]
  \item (Kasus kesamaan) Andaikan $\sup_{\intcc{-1}{1}}\abs P =
        2^{1-n}$ dengan $P$ monik berderajat $n$, dan misalkan $D = Q_n
        - P \neq 0$. Tunjukkan $(-1)^kD(\eta_k) \geq 0$ untuk setiap $k$;
        tunjukkan bahwa masing-masing dari $n$ interval
        $\intcc{\eta_{k}}{\eta_{k-1}}$ memuat sebuah nol $D$,
        dan bahwa nol yang dibagi dua interval berurutan merupakan
        titik \emph{dalam} $\eta_k$ yang $D' = 0$ pula di sana.
        Simpulkan bahwa $D$ punya $n$ nol dihitung dengan
        multiplisitasnya, jadi $D = 0$: pemberi minimumnya tepat
        $Q_n$ --- itulah \emph{teorema ekstremal Chebyshev}.
  \item Nyatakan ulang teoremanya sebagai jarak: pada $\intcc{-1}{1}$,
        \[
        d_\infty\bigl(X^n,\ \R_{n-1}[X]\bigr) = 2^{1-n},
        \]
        dengan hampiran terbaik tunggal $X^n - Q_n$; lalu tunjukkan
        lewat penyulihan afin $x = \frac{1+t}2$ bahwa pada
        $\intcc{0}{1}$ jaraknya menjadi $2^{1-2n}$.
  \item (Simpul interpolasi yang optimal) Untuk $n$ simpul $x_1, \dots,
        x_n \in \intcc{-1}{1}$, polinomial simpulnya
        $\omega(x) = \prod_i(x - x_i)$ bersifat monik berderajat $n$.
        Turunkan dari pertanyaan 12 pilihan simpul mana yang
        meminimumkan $\sup_{\intcc{-1}{1}}\abs\omega$, yakni faktor
        yang bergantung simpul pada batas galat interpolasi klasik, lalu
        berikan nilai minimumnya.
  \item Periksalah kasus $n = 2$ teoremanya dengan tangan (carilah
        $\inf_c \sup_{\intcc{-1}{1}}\abs{x^2 - c}$ secara langsung),
        lalu hitunglah secara numerik jarak pada pertanyaan 13 di
        $\intcc{0}{1}$ untuk $n = 10$. Apa yang dikatakan besarnya
        tentang grafik $x^{10}$?
\end{enumerate}

\medskip
\noindent\textbf{Bagian IV --- Bola satuan
$C(\intcc{0}{1})$.}
\begin{enumerate}[resume]
  \item Misalkan $g_k(x) = x^{2^k}$. Tunjukkan $\norm{g_k}_\infty = 1$ dan
        $\norm{g_k - g_j}_\infty \geq \frac14$ untuk $j > k$
        \emph{(nilailah di titik tempat $x^{2^k} =
        \frac12$)}: yakni sebuah barisan terbatas yang gamblang tanpa
        barisan bagian yang konvergen --- jadi bola satuan tertutupnya
        tidak \omterm{def:b2:metric:compact}{kompak}, dengan tangan kosong.
  \item (Lema Riesz, dipertajam) Misalkan $F$ subruang sejati
        \emph{berdimensi hingga} sebuah ruang bernorma
        $E$. Dengan memakai pertanyaan 1, hasilkan vektor satuan $x$
        dengan $d(x, F) = 1$ persis --- bukan sekadar $\geq 1 -
        \varepsilon$ seperti pada lema \cref{thm:b2:nvs:riesz}.
  \item Turunkan: pada setiap ruang bernorma berdimensi tak hingga ada
        barisan vektor satuan yang jarak antaranya
        $\geq 1$, lalu turunkan kembali teorema Riesz dari situ.
  \item Di $C(\intcc{0}{1})$, tunjukkan gugusan semacam itu
        secara gamblang: yakni fungsi tenda $h_n$ yang bertumpu pada
        $\bigl[\frac1{n+1}, \frac1n\bigr]$ dengan puncak bernilai $1$.
        Periksalah $\norm{h_n} = 1$, $\norm{h_n - h_m} = 1$ untuk $n
        \neq m$, lalu perhatikan bahwa $h_n \to 0$ titik demi titik
        tetapi tidak secara seragam.
  \item (Keterbatasan total gagal) Tunjukkan bahwa bola satuan tertutup
        $C(\intcc{0}{1})$ tak dapat diselimuti oleh berhingga banyak
        bola berjari-jari $\frac13$ \emph{(sebab tiap bola semacam itu
        memuat paling banyak satu $h_n$)} --- dan bandingkan dengan
        langkah keterbatasan total pada bukti
        \cref{thm:b2:metric:borellebesgue}.
\end{enumerate}

\medskip
\noindent\textbf{Bagian V --- \omterm{def:b2:nvs:norm}{Norma} yang bekerja pada matriks, dan
rangkuman.}
\begin{enumerate}[resume]
  \item Buktikan bahwa setiap \omterm{def:b2:reduction:eigen}{nilai eigen} $\lambda$ pada $A \in
        \mathcal{M}_n(\C)$ memenuhi $\abs\lambda \leq
        \vertiii A$ untuk setiap \omterm{thm:b2:nvs:continuouslinear}{norma operator}; lalu terapkan
        \cref{exo:b2:nvs:4} untuk membatasi \omterm{def:b2:reduction:eigen}{nilai eigen}
        $\begin{psmallmatrix}1 & -2\\ 3 & 1\end{psmallmatrix}$
        dan bandingkan dengan modulus sejatinya.
  \item (\omterm{def:b2:nvs:norm}{Norma} yang diselaraskan) Misalkan $A$ \omterm{def:b2:reduction:diag}{dapat didiagonalkan}, $A =
        P\,\mathrm{diag}(\lambda_1, \dots, \lambda_n)\,P^{-1}$.
        Tunjukkan bahwa $N_P(x) = \norm{P^{-1}x}_\infty$ merupakan \omterm{def:b2:nvs:norm}{norma}
        yang \omterm{thm:b2:nvs:continuouslinear}{norma operatornya} memenuhi $\vertiii A_{N_P} =
        \max_i\abs{\lambda_i}$.
  \item Turunkan: untuk $A$ yang \omterm{def:b2:reduction:diag}{dapat didiagonalkan}, $A^k \to 0$ bila
        dan hanya bila semua \omterm{def:b2:reduction:eigen}{nilai eigennya} memenuhi $\abs{\lambda_i} <
        1$ --- dan \omterm{thm:b2:nvs:finitedim}{kesetaraan norma} membuat kesimpulan itu tak
        bergantung pada \omterm{def:b2:nvs:norm}{normanya}. Periksalah pada $A =
        \frac14\begin{psmallmatrix}1 & 2\\ 2 &
        1\end{psmallmatrix}$.
  \item (Konstanta kesetaraannya meledak) Pada $\R_n[X]$, bandingkan
        $N_c(P) = \max_k \abs{a_k}$ (atas koefisiennya) dengan
        $\norm{P}_{\intcc{0}{1}}$: keduanya \omterm{def:b2:nvs:norm}{norma}, jadi keduanya
        \omterm{def:b2:nvs:equivalent}{setara} untuk tiap $n$ yang tetap; tetapi tunjukkan, dengan
        memakai pemberi minimum monik pada pertanyaan 13 di
        $\intcc{0}{1}$, bahwa konstanta terbaik $C_n$ pada $N_c \leq
        C_n\norm\cdot_{\intcc{0}{1}}$ memenuhi $C_n \geq
        2^{2n-1}$. Simpulkan dalam satu kalimat mengapa ``semua \omterm{def:b2:nvs:equivalent}{norma
        setara}'' mati dalam dimensi tak hingga.
  \item (Rangkuman) Satu kalimat untuk masing-masing: di mana kekompakan
        bola berdimensi hingga bekerja (pertanyaan 1, 12); apa yang
        dikuasai kecembungan tegas; apa yang dihancurkan gugusan pada
        pertanyaan 18--19; dan bagaimana pertanyaan 24
        mengukur kegagalannya. Sebutkan puncaknya (teorema
        ekstremal Chebyshev) lalu nyatakan di mana hampiran terbaik
        menemukan rumah modernnya (teorema proyeksi pada ruang
        Hilbert, jilid Tahun ke-3, tempat kelengkapan menggantikan
        kekompakan).

\end{enumerate}
\end{problem}
